%%%%%%%%%%%%%%%%%%%%%%%%%%%%%%%%%%%%%%%%%%%%%%%%%%%%%%%%%%%%%%%%%%%%%%%%%%%%%%%
\section{Status of Research}
\label{sec:status_of_research}

The classical van der Pauw method is a four-terminal technique for determining
the sheet resistance of thin conducting films without using a Hall-bar geometry.
Its strength lies in the separation of current injection and voltage sensing.
Ideally, the voltage contacts draw no current, so the voltage losses at the
current-injecting contacts do not enter directly into the measured sheet
resistance. This makes the method especially useful for thin films in which the
contact resistance is large, difficult to determine, or strongly dependent on
the measurement conditions.

Such contact effects are a central issue in organic field-effect transistors.
At metal--organic interfaces, charge injection is influenced by energetic
alignment, interfacial states, local morphology, and gate bias. Conventional
two-terminal transfer curves can therefore contain both channel and contact
contributions \cite{Natali2012ChargeInjection}. Bittle et al. showed that
gate-dependent contact resistance can lead to a strong overestimation of the
field-effect mobility when standard transistor formulas are applied directly to
two-terminal data \cite{Bittle2016MobilityOverestimation}. A high apparent
mobility obtained from a transfer curve is therefore not sufficient evidence for
fast transport in the semiconductor film.

The gated van der Pauw method extends the same four-terminal idea to a gated
accumulation layer. Rolin et al. introduced the approach for organic
semiconductor thin films and showed that the gate dependence of the sheet
conductance can be used to extract a long-range accumulation-layer mobility
with reduced influence from the injecting contacts \cite{Rolin2017GVDP}. In
this geometry, the gate controls the carrier density in the film, the imposed
current defines the transport path, and the transverse van der Pauw voltage
gives access to the sheet conductance of the probed region.

Later studies have broadened this idea beyond the original material system.
Bonafè et al. adapted gated van der Pauw concepts to electrolyte-gated organic
mixed ionic--electronic conductors \cite{Bonafe2021EgVDP}. Lee et al. applied
related methods to thin-film transistors for which contact resistance becomes
more important as device dimensions are reduced \cite{Lee2024GVDPIGZO}. A
subsequent voltage-driven implementation showed that channel and contact
contributions can be evaluated separately within one device
\cite{Lee2025VoltageDrivenGVDP}. Together, these studies underline the same
point: reliable mobility extraction requires a measurement strategy that can
distinguish sheet transport from contact-related voltage losses.

For a practical organic thin-film device, this separation is only meaningful if
the measurement conditions are well controlled. The prototype setup must provide
stable current injection, reliable voltage sensing, and sufficiently low
leakage. The applied current has to be large enough to produce a measurable
transverse voltage, but small enough to keep the conducting sheet close to a
linear-response regime. Gate leakage must remain small compared with the channel
current, and the extracted mobility should be reproducible for different current
levels. If the conductance curves depend strongly on the selected current, the
contact configuration, or the sweep history, the assumptions of a simple gated
van der Pauw extraction are no longer fully satisfied. For this reason, the
present thesis combines setup characterization, device fabrication,
conventional current--voltage measurements, current-driven gated van der Pauw
measurements, and additional transient/circuit tests.

